\section{Introduction}
\label{sec:introduction}

Financial mentorship is one of the strongest predictors of small-business survival, yet it remains one of the least accessible resources for young entrepreneurs across much of Africa. Formal mentorship programs are concentrated in a handful of urban hubs, typically require an existing professional network to access, and are rarely available in the languages, financial contexts, or business categories most relevant to informal-sector entrepreneurs -- market traders, smallholder agribusiness owners, and micro-manufacturers who make up a substantial share of the continent's economic activity. At the same time, smartphone penetration across the region continues to rise even where broadband infrastructure and formal financial education do not, creating an opening for AI-mediated mentorship to reach entrepreneurs that traditional programs cannot.

Large language models (LLMs) are an increasingly plausible vehicle for this kind of mentorship at scale, but a generic conversational assistant is a poor substitute for a mentor. Mentorship is not simply information delivery; it is guidance filtered through a specific, credible point of view -- the way a mentor who has personally built something from constrained circumstances reasons about risk, setbacks, and resource allocation is qualitatively different from, and more persuasive than, an undifferentiated helpful-assistant voice. This motivates the central design question behind AfriMentor AI: \emph{can a language model be constrained to consistently reason in the style of a specific, real archetype of achievement, in a way that is both measurably faithful to that archetype and genuinely useful to the person receiving guidance from it?}

We approach this question through two complementary research threads, which this paper reports on jointly. The first (Proposal 1) is methodological: we investigate whether reinforcement learning from human feedback (RLHF), conditioned on structured personality-trait profiles derived from documented African achievers, can produce a language model whose reasoning style remains consistent across a full mentorship conversation -- not just on a single prompt -- and whether such a model's outputs are distinguishable, in blind evaluation, from a generic same-capability baseline. This connects to a growing body of work on RLHF as a general alignment technique \citep{bai2022helpful,chaudhari2024rlhf} and on measuring whether personality-like traits in LLM outputs are behaviorally consistent rather than merely a consistent \emph{self-report} that fails to predict actual behavior \citep{lee2024trait,han2025illusion,jiang2023personallm}, but applies both specifically to the constraint of fidelity to a \emph{real, individual} archetype rather than a synthetic or aggregate persona -- closer in spirit to recent multi-turn reinforcement learning approaches that measurably reduce persona drift across a full conversation \citep{abdulhai2025personas} than to single-turn role-play.

The second thread (Proposal 2) is empirical and user-centered: what do African micro-entrepreneurs, specifically those with limited formal education and constrained access to data and devices, actually want from an AI mentorship product, and does a personality-constrained persona change engagement or perceived usefulness relative to a generic assistant? This question sits within a broader literature on LLM-based tutoring and educational support systems, much of which has been developed and evaluated through conversation-based tutoring architectures with explicit student modeling and retrieval-augmented content \citep{park2024tutoring,liu2025lpitutor}, and through systematic reviews highlighting both the promise and the open ethical and accessibility questions of LLM-based personalized learning \citep{sharma2025review}. We are specifically interested in whether those findings transfer to an informal-economy, low-bandwidth, voice-first context, or whether they need to be substantially rethought.

The system built to support this research, AfriMentor AI, is a smartphone-first mentorship platform with three defining constraints, chosen specifically to make both research threads testable rather than purely aspirational: (1) a knowledge base restricted to African-authored financial and entrepreneurship content, so that guidance is contextually grounded rather than generically Western; (2) a voice-note-first interaction model, chosen to match the primary user's actual device and connectivity constraints rather than assuming typed, always-online interaction; and (3) an explicit, inspectable mapping from a small set of achiever archetypes to persona-conditioning parameters, so that persona fidelity is something we can actually measure rather than only claim.

The remainder of this paper is organized as follows. Section~\ref{sec:related-work} situates this work relative to prior research on RLHF alignment, personality modeling in LLMs, generative agents, and LLM-based tutoring. Section~\ref{sec:method} (forthcoming) describes the persona-conditioning and RLHF methodology in detail. Section~\ref{sec:evaluation} (forthcoming) reports the persona-fidelity and user-study evaluation. We close with a discussion of limitations specific to deploying this kind of system in low-resource, informal-economy contexts.
